\section{Usage Notes}
\label{sec:usage_notes}

\subsection{Recommended Uses}

This dataset is prepared for secondary analysis and does not present new biological findings.
Documented intended reuse includes methods development, multimodal MRI analysis, quality-control benchmarking, and reproducible neuroimaging workflows that start from BIDS inputs.
The release includes structural, functional, diffusion, and field-map MRI organized according to BIDS (BIDSVersion 1.9.0), with a longitudinal design of up to two sessions (\texttt{ses-01}, \texttt{ses-02}) when both visits exist.

\subsection{Important Considerations}

Users should account for the following documented constraints when planning analyses.

\begin{itemize}
    \item \textbf{Longitudinal coverage is incomplete.}
    Of 84 subjects in the BIDS tree, 51 have both sessions, 32 have \texttt{ses-01} only, and 1 has \texttt{ses-02} only.
    Missing sessions reflect incomplete participant follow-up, not omitted files.
    BIDS Validator warnings such as \texttt{MISSING\_SESSION} and \texttt{INCONSISTENT\_SUBJECTS} are expected.

    \item \textbf{Task events are not available for every functional run.}
    For \texttt{task-fmri} (laboratory visual grating / checkerboard paradigm), event timing files accompany 514 of 553 magnitude BOLD series (92.9\%); runs without an unambiguous protocol-to-acquisition mapping intentionally lack \texttt{events.tsv}.
    For \texttt{task-movie}, design-level \texttt{events.tsv} files are provided for all magnitude movie BOLD runs, but onsets follow a documented run-level / continuous-movie timing policy and are \textbf{not} millisecond TTL-locked onsets.
    Movie stimulus video files are \textbf{not} redistributed (third-party copyright); protocol code and timing metadata are included under \texttt{code/task-movie/}.
    For \texttt{task-rest}, no \texttt{events.tsv} are released by design (continuous fixation without discrete trials); validator \texttt{EVENTS\_TSV\_MISSING} warnings are expected.
    Control runs lack recoverable trial timing and intentionally have no \texttt{events.tsv}.
    Incomplete task coverage includes \texttt{sub-055} (no rest) and \texttt{sub-063} (no grating), as documented in the release README.

    \item \textbf{Public package scope.}
    The public package is intended as BIDS data, with defaced anatomical derivatives for sharing where applicable, and \textbf{must not} redistribute identifiable source DICOM (for example under \texttt{raw\_original/}).
    Defaced structural volumes are generated under a derivatives tree; original BIDS anatomicals are not overwritten by defacing.

    \item \textbf{Derivatives versus raw acquisitions.}
    Preprocessed diffusion volumes under \texttt{derivatives/tractoflow\_normalize\_dwi/} were generated by an upstream workflow, are present for only a partial subset of subjects, and should be treated as derivatives---not as replacements for raw BIDS \texttt{dwi/} runs.
    Upstream analysis products such as HCP structural preprocessing outputs, Benson's Atlas labels, tractography, or connectivity matrices should not be assumed to be present unless listed as verified contents of the release tree.

    \item \textbf{Quality control.}
    Automated QC outputs (including BIDS validation, MRIQC where available, Pizarro T1w screening, DWI QC, defacing audit, and privacy audit) are used to describe and prioritize inspection.
    They are \textbf{not} automatic exclusion criteria for this release description.
    MRIQC and cohort-wide DWI QC summaries were not finished at the time of the current release README.
    Residual visual review of face removal remains recommended. A prior paired T1w audit covered 356/356 volumes then in that campaign; the live packaged tree has 382 T1w after five visual-QC withdrawals. Selected anatomicals with insufficient face removal were re-run with \texttt{pydeface} on 2026-08-26 (28 volumes; \texttt{docs/quality\_control/defacing\_redo\_2026-08-26.md}).

    \item \textbf{De-identification.}
    De-identification follows a DICOM PS3.15-oriented custom workflow and is \textbf{not} a claim of certified full compliance with the DICOM PS3.15 Basic Profile.
    Siemens PhysioLog objects that passed documented conversion gates are distributed as BIDS physiology sidecars; peripheral PMU logs and failed-gate PhysioLogs remain source-only.

    \item \textbf{License / access.}
    This dataset is intended for contractual release via FITBIR.
    Access is governed by FITBIR data-access agreements and study-specific sharing terms (see the dataset \texttt{LICENSE}).
    It is \textbf{not} released under CC0.
    Repository accession / how-to-cite strings remain to be inserted when the FITBIR deposit is finalized.
\end{itemize}

\subsection{Software and Tools}

Release documentation records the following software used to convert, organize, validate, or quality-check the dataset (versions as stated in the release README and related reports):

\begin{itemize}
    \item \texttt{neuro\_pipeline} (dataset\_description GeneratedBy version 2.1.0) for DICOM-to-BIDS organization;
    \item dcm2niix for NIfTI + JSON conversion (sidecar provenance typically records \texttt{v1.0.20230411});
    \item bids-validator 1.15.0;
    \item MRIQC (approximately version 24.x; coverage partial at the time of documentation);
    \item pydeface 2.0.0 (Compute Canada build) for anatomical defacing derivatives;
    \item Pizarro T1w screening with ONNX Runtime 1.23.2;
    \item MRtrix3 tools for DWI QC;
    \item Python 3.11.x in documented report environments.
\end{itemize}

Release-facing paradigm code and dictionaries live under \texttt{release\_dataset/code/} (including \texttt{code/task-grating/}, \texttt{code/task-movie/}, and related condition dictionaries).
Pipeline, cleaning, QC, and audit scripts are documented in the project workspace (examples listed in the release README under Code availability).
Users working from BIDS inputs may also use standard BIDS-compatible neuroimaging tools of their choice; the items above are those recorded for the construction and validation of this release, not a mandatory analysis stack.
